% Ad Atticum 4.16 --- headnote
% ad-atticum-04-16, late June or early July 54 BC, Rome

\textit{Cicero to Atticus, written from Rome in late June or early
July 54 BC --- the Perseus dateline gives ``end of June or in
Quintilis 700 (54).'' Atticus, having left Rome later than planned
(as 4.14 already complained), has now sailed across the Adriatic
and through Epirus, sent a flurry of short notes back from
Buthrotum, and is on his way further east into Asia. Cicero opens
by signaling his own pressure of business --- the letter is dictated
to a secretary --- and works methodically through the substantive
letter Paccius brought him: the introduction of Paccius into
Cicero's intimate circle on Atticus's word, and then the long
question Atticus has been pressing about \textit{De Re Publica}.}

\textit{Sections 2 and 3 are the rarest thing in the correspondence:
Cicero on his own compositional method, mid-draft. He explains why
Varro cannot simply be inserted as a speaker into a dialogue whose
dramatic date is fixed in the previous century (the cast is
Africanus, Philus, Laelius, Manilius, with the younger Tubero,
Rutilius, Scaevola, and Fannius), and floats the alternative of a
dedicatory proem in the manner of Aristotle's \textit{exoteric}
works. He then defends, against Atticus's objection, the
disappearance of the elder Scaevola after Book 1 by direct appeal
to Plato's handling of Cephalus in the \textit{Republic} ---
``our divine Plato.'' From section 4 the letter turns to the
familiar Roman bulletin: legal news about C. Cato's serial
acquittals, the bringing of Drusus and Scaurus to trial, the
shaky senatus consultum on provinces, and the deadlocked consular
field --- Messalla, Scaurus, Domitius, Memmius --- in which Caesar's
soldiers, Pompey's Gaul, and the threat of pushing the elections
into Caesar's return all already figure as ordinary political
currency. The letter closes with a gentle but real scolding for
Atticus's decision to extend the journey to Asia rather than
sending agents in his place.}
